% Generated draft: requires mathematical and source review.
\documentclass[11pt]{article}
\usepackage[margin=25mm]{geometry}
\usepackage{fontspec}
\setmainfont{Latin Modern Roman}
\usepackage{amsmath,amssymb,mathtools,unicode-math}
\setmathfont{Latin Modern Math}
\usepackage{enumitem,needspace,xurl,hyperref}
\setlength{\parindent}{0pt}
\setlength{\parskip}{5pt}
\setlength{\emergencystretch}{3em}
\newcommand{\N}{\mathbb{N}}
\newcommand{\Z}{\mathbb{Z}}
\newcommand{\Q}{\mathbb{Q}}
\newcommand{\R}{\mathbb{R}}
\newcommand{\C}{\mathbb{C}}
\newcommand{\F}{\mathbb{F}}
\newcommand{\A}{\mathbb{A}}
\newcommand{\PP}{\mathbb{P}}
\newcommand{\E}{\mathbb{E}}
\newcommand{\eps}{\varepsilon}
\newcommand{\dd}{\,\mathrm{d}}
\begin{document}
\section*{UnsolvedMath ID: 700967}
\textbf{Source problem code:} catalogue entry
\subsection{Definitions and mathematical statement}
Let \(G\) be a group equipped with a topology for which multiplication \(G\times G\to G\) and inversion \(G\to G\) are continuous, and let \(e\) denote its identity. The source does not specify any separation axiom; none is added in the formulation below, and whether a Hausdorff convention is intended remains unresolved.

An ultrafilter on the underlying set of \(G\) is a family \(\mathcal U\subseteq\mathcal P(G)\) such that \(G\in\mathcal U\), \(\varnothing\notin\mathcal U\), the family is closed under finite intersections and passage to supersets, and for every \(S\subseteq G\), either \(S\in\mathcal U\) or \(G\setminus S\in\mathcal U\). It is free if it contains no finite subset of \(G\). It converges to \(e\) if every neighborhood of \(e\) belongs to \(\mathcal U\). Define
\[
\mathscr U_e(G)=\{\mathcal U: \mathcal U\text{ is a free ultrafilter on }G\text{ converging to }e\}.
\]
A subset \(D\subseteq G\) is dense if its closure in \(G\) equals \(G\).

\textbf{Question (Kourovka Notebook 13.49).} Is it true that, for every topological group \(G\), if \(\mathscr U_e(G)\) is infinite, then there exist subsets \(A,B\subseteq G\) satisfying
\[
G=A\cup B,\qquad A\cap B=\varnothing,\qquad \overline A=\overline B=G?
\]
Here infinitude refers to infinitely many distinct ultrafilters on the underlying set, and the partition is not required to have any additional regularity.
\subsection{Short English statement}
Must every topological group with infinitely many distinct free ultrafilters converging to its identity admit a partition into two dense subsets? The source leaves the separation convention for topological groups unspecified.
\subsection{Sources}
\begin{itemize}
\item Local UnsolvedMath registry, supplied complete source record: permanent ID \(700967\), exact source code \texttt{catalogue entry}, title \emph{Kourovka Notebook Problem 13.49}.
\item Supplied reference: \emph{The Kourovka Notebook (Unsolved Problems in Group Theory)}, 13th issue (1995), Problem 13.49; also identified in the record with the 21st issue edition (2026), \texttt{arXiv:1401.0300v48}. The collection description names E. I. Khukhro and V. D. Mazurov as editors.
\item Attribution is retained as supplied: the statement begins with ``(V. I. Malykhin)'', while the background and proposer field identify I. V. Protasov as proposer, with proposed year 1995. The relationship between these attributions is not explained.
\item The registry labels the problem \emph{open}. This is a supplied status label, not an independently verified assessment. The cited edition and source contents have not been independently checked for this draft.
\end{itemize}
\end{document}
